% ============================================================
%  "Señales visuales, estimación de demanda y diseño de
%   mecanismos en la asignación escolar de Bogotá"
%  Bautista, Fuentes, Melo (2026)
% ============================================================

\documentclass[12pt,a4paper]{article}

\usepackage[utf8]{inputenc}
\usepackage[T1]{fontenc}
\usepackage[main=spanish]{babel}
\usepackage[round,authoryear]{natbib}
\usepackage{hyperref}
\usepackage{microtype}
\usepackage{booktabs}
\usepackage{amsmath}
\usepackage{geometry}
\geometry{margin=2.5cm}
\setlength{\parskip}{4pt}
\usepackage{lmodern}

\begin{document}

\title{\textbf{Señales visuales, estimación de demanda y diseño de mecanismos en la asignación de cupos escolares de Bogotá}}
\author{Bautista, D. \and Fuentes, J. \and Melo, J.}
\date{2026}
\maketitle

% ============================================================
\section*{1.\quad Introducción}
% ============================================================

El sistema de educación pública de Bogotá agrupa 413~instituciones en 740~sedes
distribuidas en 20~localidades \citep{SED2023caracterizacion}. La demanda oficial
muestra contracción sostenida ($-1{,}2$\,\% anual 2019--2023), pero las localidades
de menor ingreso --- Kennedy, Bosa, Ciudad~Bolívar --- concentran más del 40\,\%
de los cupos nuevos solicitados.

La asignación de cupos está regulada por la \citet{Resolucion1587} mediante un
índice de priorización que combina condición socioeconómica (SISBEN) y proximidad
geográfica, formalizable como DA \citep{GaleShapley1962} con prioridad
lexicográfica. Sin embargo, el mecanismo opera sobre preferencias que pueden
estar sistemáticamente sesgadas: la literatura muestra que las familias de
menores ingresos son más sensibles a señales no académicas al elegir colegio
\citep{Hastings2009}, y la apariencia física del establecimiento se convierte
en el \emph{proxy} de calidad más accesible cuando la información académica es
costosa \citep{Naik2017,Dubey2016}. Si las familias más vulnerables eligen por
apariencia en lugar de calidad, el sesgo perjudica exactamente a quienes el
mecanismo intenta proteger \citep{HastingsWeinstein2008}.

Este trabajo aborda dos preguntas: (i)~¿en qué medida la apariencia visual de
los colegios distorsiona las preferencias?, y (ii)~¿puede un mecanismo de
asignación corregir esa distorsión? Construimos señales visuales a partir de
imágenes GSV procesadas con CLIP~\citep{Radford2021} y segmentación semántica,
estimamos su efecto sobre la demanda mediante un modelo BLP
\citep{Berry1994,ConlonGortmaker2020}, y evaluamos mecanismos estándar --- BM,
DA y SED-lex --- junto con una regla aprendida de tipo WP sobre 99{.}890~familias
de primer ingreso calibradas con la Encuesta Multipropósito~2021 \citep{DANE2021EM}.

Los resultados muestran que la seguridad percibida del colegio predice
significativamente la demanda ($\hat{\beta}\!=\!+0{,}11$, $p\!<\!0{,}01$),
pero ningún mecanismo estándar corrige directamente el canal visual. Proponemos
entonces una WP-Rule estable cuya prioridad de colegios combina ingreso y distancia,
e incorpora una corrección visual calibrada para neutralizar la correlación entre
ingreso y seguridad percibida asignada.


% ============================================================
\section*{2.\quad Datos y metodología}
% ============================================================

\subsection*{2.1\quad Datos de colegios y demanda}

La unidad de análisis es el \emph{establecimiento educativo} (no la sede). El
directorio SED~2024 reporta 2{.}242~sedes agrupadas en establecimientos; al filtrar
el sector oficial y agregar sedes a establecimiento~--- tomando la sede principal como
punto geográfico representativo~---, se obtienen inicialmente~408~establecimientos
tras fusionar matrícula, demanda, puntajes Saber~11
(promedio 2020--2022--2023) y controles socioeconómicos de la Encuesta
Multipropósito~2021 por UPZ. Se excluyen 23~colegios rurales (Sumapaz y sedes
aisladas sin cobertura EM2021) y 3~establecimientos sin cobertura GSV,
resultando en \textbf{382~colegios oficiales urbanos} como muestra final.
La exclusión de los rurales es sustantiva: un análisis preliminar con NMF sobre embeddings VGG19~\citep{Simonyan2015}
mostró que la totalidad de la señal visual capturada correspondía al contraste
rural/urbano, no a variación dentro del entorno urbano.

Para la estimación estructural, se construye la cuota de
mercado $s_j = D_j / M_t$ por localidad y se aplica la inversión de
\citet{Berry1994}: $\delta_j = \ln s_j - \ln s_0$, donde $s_0$ es la
outside option estimada por localidad. Los controles incluyen puntaje
Saber~11, homicidios~(log), distancia a SITP~(log), competencia privada y
condición técnica.

\subsection*{2.2\quad Pipeline de visión computacional}

Se descargaron imágenes GSV de fachada para cada colegio oficial
(10~headings por sede a~$0^\circ, 36^\circ,\ldots,324^\circ$;
640$\times$640~px, FOV\,=\,$90^\circ$). Para cada establecimiento se procesan las 10~imágenes disponibles y se
agregan los scores a nivel de establecimiento. Los~3~establecimientos sin GSV reciben imputación por vecindad en radio de~2~km.

Se emplean dos estrategias complementarias de extracción visual:

\begin{enumerate}
  \item \textbf{CLIP ViT-B/32} \citep{Radford2021}: se computan 4~scores de
        percepción por establecimiento mediante prompts textuales calibrados
        (mantenimiento, vegetación percibida, modernidad y seguridad percibida).
        Cada score mide la similitud coseno entre el embedding de la imagen y el
        del prompt; a nivel de establecimiento se conserva el máximo observado
        entre las imágenes disponibles.
  \item \textbf{Segmentación semántica (DeepLabV3+/Cityscapes):} se extraen
        proporciones de píxeles de cada clase semántica por imagen. Tras análisis
        de colinealidad con las variables CLIP, se retienen 2~proporciones
        informativas: infraestructura vial y cerramiento.
\end{enumerate}


\subsection*{2.3\quad Estimación de demanda}

\paragraph{Berry OLS (segunda etapa).}
La utilidad media $\delta_j$ obtenida por inversión de Berry se descompone
mediante OLS en cuatro especificaciones que añaden sucesivamente las señales
visuales a los controles:
\begin{equation}
\delta_j \;=\; \beta_0 + \beta_{\text{seg}}\,\text{seg}_j
  + \beta_{\text{veg}}\,\text{veg}_j
  + \beta_{\text{vial}}\,\text{vial}_j
  + \beta_{\text{cerr}}\,\text{cerr}_j
  + \beta_q\, q_j + \mathbf{B}_C\,\mathbf{c}_j + \xi_j
\label{eq:berry_ols}
\end{equation}
donde $\text{seg}_j$ y $\text{veg}_j$ son scores CLIP de seguridad y vegetación
percibidas; $\text{vial}_j$ y $\text{cerr}_j$ son proporciones Cityscapes de
infraestructura vial y cerramiento; $q_j$ mide calidad académica; y
$\mathbf{c}_j$ agrupa controles de entorno escolar. Los resultados se presentan
en la Tabla~\ref{tab:berry_blp}.

\paragraph{IV-BLP (primera etapa).}
Para capturar la heterogeneidad en preferencias por ingreso, se estima un modelo
BLP con micro-momentos \citep{Berry1994,ConlonGortmaker2020}:
\begin{equation}
U_{ij} = \delta_j + \pi_1\, y_i\, \text{seg}_j
  + \lambda_0\,\log(1{+}d_{ij}) + \lambda_1\, y_i\,\log(1{+}d_{ij})
  + \varepsilon_{ij}
\label{eq:blp}
\end{equation}
donde $\delta_j$ se recupera por contraction mapping y los parámetros no lineales
$(\pi_1, \lambda_0, \lambda_1)$ se estiman por GMM. La calidad $q_j$ se instrumenta
con la calidad media de rivales ponderada por $1/d$ dentro de la localidad.

\subsection*{2.4\quad Familias y utilidades}

La EM2021 registra hogares con hijos en colegio oficial en Bogotá; al expandir
por~$\text{FEX}_C$ y conservar familias que buscan cupo de primer ingreso se
obtienen 99{.}890~familias representativas, ubicadas en manzanas coherentes con
su estrato y UPZ. A cada familia se le
asigna un grupo SISBEN ($A$--$D$) a partir del ingreso per~cápita. La utilidad
sigue la ecuación~\eqref{eq:blp} con parámetros del IV-BLP; las preferencias
son el ranking de colegios accesibles por utilidad ($d_{ij}$~Haversine). La
capacidad de cada colegio es $\lfloor M_j/13\rceil$ (mínimo~5). En la
evaluación de mecanismos, el universo efectivo está dado por los colegios que
aparecen en las preferencias top-20 y tienen señal visual observada: 380
colegios y 113{.}857~cupos.

\subsection*{2.5\quad Mecanismos de asignación}

Evaluamos cuatro mecanismos. Los tres primeros son estándar en la literatura de
\emph{school choice}; el cuarto es una contribución de este trabajo.

\paragraph{Boston (BM).}
Asignación inmediata con aceptaciones definitivas y prioridad geográfica.
No es incentivo-compatible (IC): las familias pueden beneficiarse al no reportar
sus preferencias verdaderas \citep{AbdulSonmez2003}.

\paragraph{Deferred Acceptance (DA).}
Gale-Shapley estudiantil con aceptaciones provisionales y prioridad geográfica.
IC, estable (cero blocking pairs) y óptimo entre matchings estables para los
estudiantes \citep{AbdulSonmez2003}.

\paragraph{SED-lex.}
DA con prioridad lexicográfica SISBEN ($A>B>C>D$) y desempate por proximidad
geográfica. Réplica del mecanismo vigente de la SED \citep{Resolucion1587}.
Hereda las propiedades de DA (IC, estabilidad) y prioriza familias vulnerables.

\paragraph{Weighted Polytope Rule (WP).}
La regla WP sigue el marco de \citet{NarasimhanAgarwalParkes2016}: aprende una
función de bienestar lineal y la maximiza sobre el politopo de asignaciones
estables. A diferencia de SED-lex, la prioridad de los colegios no es
lexicográfica en SISBEN. Definimos una prioridad continua que favorece menores
ingresos y menor distancia, e incorpora una corrección visual:
\begin{equation}
P^\theta_{ij}
=
-\tilde y_i
- \widetilde d_{ij}
- \theta\,\tilde y_i v_j ,
\label{eq:priority_id}
\end{equation}
donde $\tilde y_i$ es ingreso per~cápita estandarizado, $\widetilde d_{ij}$ es
distancia estandarizada, y $v_j$ es seguridad percibida. Valores mayores de
$P^\theta_{ij}$ implican mayor prioridad del hogar $i$ en el colegio $j$. El
parámetro $\theta$ se calibra en mercados sintéticos Bogotá-like para que la
correlación entre ingreso y seguridad percibida asignada sea cercana a cero.

Dada esta prioridad, la WP maximiza:
\begin{equation}
f_{\mathrm{WP}}(\succ;W)
\in
\arg\max_{x \in \widehat{\Omega}^{\mathrm{st}}(\succ,P^\theta)}
\sum_{ij} \lambda_{ij}(W)\,x_{ij}.
\label{eq:wp_rule}
\end{equation}
Usamos una parametrización parsimoniosa:
\begin{equation}
\lambda_{ij}(W)
=
a R_{ij}
+ b P^\theta_{ij}
+ d_v \tilde y_i v_j .
\label{eq:wp}
\end{equation}
Aquí $R_{ij}$ aumenta cuando la familia rankea mejor al colegio, $P^\theta_{ij}$ mide la
prioridad ingreso--distancia ajustada por \eqref{eq:priority_id}, y el término
$d_v\tilde y_i v_j$ permite que el objetivo WP penalice asignaciones que vuelvan
a correlacionar ingreso y seguridad percibida.

El StructSVM aprende $W=(a,b,d_v)$ a partir de mercados sintéticos. La estabilidad
no se relaja: queda garantizada por las restricciones de
$\widehat{\Omega}^{\mathrm{st}}(\succ,P^\theta)$, que incorporan cupos y ausencia
de pares bloqueantes bajo la prioridad ingreso--distancia ajustada
\citep{VandeVate1989}. Como benchmarks incluimos WP\textsubscript{baseline}, sin
corrección visual en el objetivo ($d_v=0$), y WP\textsubscript{equal}, con pesos
iguales sobre preferencias y prioridad.

En la aplicación a datos reales no resolvemos el LP completo de
\eqref{eq:wp_rule} sobre las 99{.}890 familias, porque el número de variables
familia--colegio y restricciones de estabilidad crece al orden de millones aun
usando listas top-20. En su lugar aplicamos la parte escalable de la regla
aprendida: corremos DA con la prioridad calibrada $P^\theta_{ij}$. Esta versión,
que llamamos WP-priority, no reentrena el modelo y preserva estabilidad exacta
bajo la prioridad ingreso--distancia--visual, pero debe interpretarse como la
implementación a escala real de la prioridad WP, no como la solución del LP
global de WP.


% ============================================================
\section*{3.\quad Resultados}
% ============================================================

\subsection*{3.1\quad Estimación de demanda}

La Tabla~\ref{tab:berry_blp} presenta los resultados de la segunda etapa
(Berry~OLS) y los parámetros heterogéneos del IV-BLP.

\begin{table}[h!]
\centering
\caption{Estimación de demanda escolar (382 colegios urbanos)}
\label{tab:berry_blp}
\small
\begin{tabular}{lcccc|c}
\toprule
 & \multicolumn{4}{c|}{\textit{Berry OLS --- $\delta_j$}} &
   \textit{IV-BLP} \\
\textbf{Variable} & M0 & M1 & M2 & M3 & $\theta$ \\
\midrule
\multicolumn{6}{l}{\textit{Señales visuales}} \\
\quad Seg.\ percibida (CLIP) & --- &
  0{,}108$^{***}$ & --- & 0{,}110$^{***}$ &  \\
\quad Veg.\ percibida (CLIP) & --- &
  $-$0{,}001 & --- & $-$0{,}002 &  \\
\quad Infra.\ vial (CS) & --- & --- &
  0{,}023 & 0{,}019 &  \\
\quad Cerramiento (CS) & --- & --- &
  0{,}023 & 0{,}030 &  \\
\addlinespace
\multicolumn{6}{l}{\textit{Calidad y controles}} \\
\quad Saber 11 ($q_j$) &
  $-$0{,}183$^{***}$ & $-$0{,}183$^{***}$ &
  $-$0{,}181$^{***}$ & $-$0{,}180$^{***}$ &
  $-$0{,}184$^{***}$ \\
\quad Homicidios (log) &
  $-$0{,}440$^{***}$ & $-$0{,}432$^{***}$ &
  $-$0{,}443$^{***}$ & $-$0{,}435$^{***}$ &
  $-$0{,}440$^{***}$ \\
\quad Técnico &
  0{,}485$^{***}$ & 0{,}441$^{***}$ &
  0{,}483$^{***}$ & 0{,}439$^{***}$ &
  0{,}445$^{***}$ \\
\quad Compet.\ privada &
  $-$0{,}443$^{***}$ & $-$0{,}420$^{***}$ &
  $-$0{,}443$^{***}$ & $-$0{,}418$^{***}$ &
  $-$0{,}417$^{***}$ \\
\addlinespace
\multicolumn{6}{l}{\textit{Parámetros heterogéneos (BLP)}} \\
\quad $\pi_1$ ($y_i \times$ seg) & & & & & $-$0{,}028 \\
\quad $\lambda_0$ (distancia) & & & & & $+$0{,}024 \\
\quad $\lambda_1$ ($y_i \times$ dist) & & & & & $-$0{,}094 \\
\midrule
R$^2_{\text{adj}}$ & 0{,}460 & 0{,}469 & 0{,}459 & 0{,}467 & --- \\
First-stage $F$ & & & & & 18{,}8 \\
\bottomrule
\end{tabular}
\begin{minipage}{\linewidth}
\vspace{4pt}\footnotesize
\textit{Nota.} Variables estandarizadas (media~0, DE~1). $^{***}$\,$p<0{,}01$.
Berry OLS: variable dependiente $\delta_j = \ln s_j - \ln s_0$.
IV-BLP: $q_j$ instrumentado con calidad media de rivales ponderada por $1/d$
en la localidad.
CS\,=\,Cityscapes.
\end{minipage}
\end{table}

Tres hallazgos centrales.
\textbf{Primero}, la seguridad percibida (CLIP) es la única señal visual
estadísticamente significativa ($+0{,}110$, $p\!<\!0{,}01$): los colegios que
lucen más seguros atraen más demanda, independientemente de su calidad académica.
\textbf{Segundo}, el puntaje Saber~11 tiene coeficiente \emph{negativo}
($-0{,}180$): controlando por zona, las familias no eligen por resultados
académicos.
\textbf{Tercero}, el IV-BLP confirma heterogeneidad por ingreso:
$\hat{\lambda}_1\!=\!-0{,}094$ indica que la distancia penaliza más a familias
de mayor ingreso, y $\hat{\pi}_1\!=\!-0{,}028$ sugiere que familias de mayor
ingreso responden marginalmente menos a la seguridad percibida. Los coeficientes
lineales son robustos a la instrumentación.

\subsection*{3.2\quad WP-Rule}

La WP-Rule se evalúa primero en mercados sintéticos \emph{Bogotá-like}, donde es
posible observar si una prioridad diseñada puede romper el canal visual sin
perder estabilidad. La prioridad ingreso--distancia--visual se calibra para que
la asignación estable no herede una correlación sistemática entre ingreso y
seguridad percibida; luego el StructSVM aprende los pesos
$W=(a,b,d_v)$ de la función objetivo.

En validación, los tres criterios WP entregan resultados prácticamente idénticos:
\[
\begin{array}{lcccc}
\toprule
 & \text{corr}(y,v) & \text{corr}(y,q) & \bar q & \text{rank} \\
\midrule
\text{WP}_{\text{learned}}  & -0{,}0083 & 0{,}0508 & -0{,}0098 & 7{,}00 \\
\text{WP}_{\text{baseline}} & -0{,}0083 & 0{,}0508 & -0{,}0098 & 7{,}00 \\
\text{WP}_{\text{equal}}    & -0{,}0080 & 0{,}0511 & -0{,}0098 & 7{,}00 \\
\bottomrule
\end{array}
\]
La correlación ingreso--seguridad queda cerca de cero sin introducir pares
bloqueantes. Además, la similitud entre WP\textsubscript{learned},
WP\textsubscript{baseline} y WP\textsubscript{equal} muestra que la corrección
proviene principalmente de la prioridad que define el conjunto estable, no de
pequeñas diferencias en los pesos finales del objetivo.

Fuera de validación, la regla mantiene el mismo patrón: en mercados Bogotá-like de
test se obtiene $\text{corr}(y,v)=0{,}0129$ y
$\text{corr}(y,q)=0{,}1228$. En el escenario diagnóstico ortogonal
($\rho=0$), donde seguridad percibida y calidad se generan sin correlación, la WP
produce $\text{corr}(y,v)=0{,}0288$ y $\text{corr}(y,q)=0{,}0729$.

La robustez a distintos valores de
$\rho=\operatorname{corr}(v_j,q_j)$ refuerza esta lectura. Para
WP\textsubscript{learned}, $\text{corr}(y,v)$ pasa de 0{,}0011 cuando $\rho=0$ a
0{,}0993 cuando $\rho=1$, siempre con cero pares bloqueantes y acceso A/B cercano
a 0{,}98--0{,}99. La conclusión de esta sección es que WP funciona como prueba de
diseño: una prioridad estable, construida con ingreso, distancia y señal visual,
puede contener el sorting visual en mercados controlados.

\subsection*{3.3\quad Comparación de mecanismos}

Comparamos BM, DA, SED-lex y WP en dos ejercicios. El primero usa mercados
sintéticos Bogotá-like; el segundo aplica los mecanismos a las preferencias y
capacidades reales. En ambos casos, las diferencias provienen de la regla de
prioridad y de si el mecanismo impone estabilidad.

\begin{table}[h!]
\centering
\caption{Comparación de mecanismos en mercados sintéticos Bogotá-like}
\label{tab:wp_comparacion}
\footnotesize
\begin{tabular}{lccccc}
\toprule
Mecanismo & $\text{corr}(y,v)$ & $\text{corr}(y,q)$ & $\bar q$ & Rank & BP \\
\midrule
BM                  & $-0{,}033$ & $-0{,}030$ & $0{,}052$ & 4{,}58 & 24{,}8 \\
DA                  & $ 0{,}002$ & $-0{,}005$ & $0{,}052$ & 8{,}07 & 0{,}0 \\
SED-lex             & $-0{,}077$ & $ 0{,}129$ & $0{,}052$ & 5{,}51 & 0{,}0 \\
WP\textsubscript{learned} & $ 0{,}018$ & $ 0{,}078$ & $0{,}052$ & 6{,}91 & 0{,}0 \\
\bottomrule
\end{tabular}
\begin{minipage}{\linewidth}
\vspace{4pt}\footnotesize
\textit{Nota.} Rank es la posición promedio del colegio asignado en la lista de
preferencias del hogar; BP reporta pares bloqueantes promedio por mercado.
\end{minipage}
\end{table}

En el escenario sintético, BM obtiene el mejor rank promedio, pero genera pares
bloqueantes. DA, SED-lex y WP son estables. SED-lex mejora el rank frente a DA,
pero aumenta la correlación ingreso--calidad; esto no contradice su prioridad por
vulnerabilidad, porque la regla respeta las preferencias declaradas. Si esas
preferencias están mediadas por seguridad percibida, distancia u oferta local,
priorizar a hogares vulnerables dentro de sus opciones elegidas no necesariamente
corrige el sorting agregado hacia calidad académica. WP\textsubscript{learned}
mantiene estabilidad y contiene la correlación ingreso--seguridad, con un costo de
rank frente a BM y SED-lex.

En datos reales, la lectura cambia. El universo efectivo contiene
99{.}890~familias de primer ingreso, 380~colegios presentes en las preferencias
top-20 y 113{.}857~cupos. Todos los mecanismos asignan 96{.}022~familias; por
tanto, la comparación relevante no es el volumen total asignado, sino el sorting
por ingreso, calidad y seguridad percibida.

\begin{table}[h!]
\centering
\caption{Comparación de mecanismos en datos reales}
\label{tab:wp_real}
\small
\begin{tabular}{lccccc}
\toprule
Mecanismo & $\text{corr}(y,v)$ & $\text{corr}(y,q)$ & $\bar q$ & Rank & BP \\
\midrule
BM                  & $0{,}0188$ & $0{,}1067$ & $256{,}67$ & 1{,}83 & 4{.}148 \\
DA                  & $0{,}0242$ & $0{,}1329$ & $256{,}68$ & 2{,}63 & 0 \\
SED-lex             & $0{,}0016$ & $0{,}0979$ & $256{,}67$ & 1{,}91 & 0 \\
WP-priority         & $0{,}0181$ & $0{,}1043$ & $256{,}67$ & 1{,}92 & 0 \\
\bottomrule
\end{tabular}
\begin{minipage}{\linewidth}
\vspace{4pt}\footnotesize
\textit{Nota.} WP-priority aplica la prioridad ingreso--distancia--visual
calibrada en la WP-Rule y corre DA sobre datos reales; no resuelve el LP global
de WP a escala completa.
\end{minipage}
\end{table}

SED-lex es el mecanismo que más neutraliza el canal visual:
$\text{corr}(y,v)=0{,}0016$. WP-priority mejora frente a DA y mantiene estabilidad
exacta, pero no supera a SED-lex. La diferencia con el ejercicio sintético es
informativa: con preferencias reales restringidas por localidad, la prioridad
lexicográfica por vulnerabilidad y el desempate geográfico ya corrigen casi por
completo el sorting ingreso--seguridad. WP-priority queda como una alternativa
estable y continua, pero el mecanismo vigente contiene una palanca redistributiva
más fuerte en este contexto.


% ============================================================
\section*{4.\quad Conclusiones y limitaciones}
% ============================================================

Este trabajo muestra que las señales visuales del entorno escolar contienen
información conductual relevante: la seguridad percibida medida con CLIP predice
sobredemanda aun controlando por calidad académica y entorno urbano
($\hat\beta_{\text{seg}}\approx 0{,}11$, $p<0{,}01$). Al mismo tiempo, el puntaje
Saber~11 no aumenta la demanda condicional. La implicación es directa: antes de
la matrícula, las familias parecen responder más a señales visibles y accesibles
que a medidas académicas difíciles de observar.

La segunda contribución es de diseño de mecanismos. En mercados sintéticos,
la WP-Rule muestra que una prioridad ingreso--distancia--visual puede reducir el
sorting ingreso--seguridad manteniendo estabilidad. Sin embargo, al aplicar los
mecanismos sobre datos reales, SED-lex es el mecanismo que mejor neutraliza el
canal visual: $\text{corr}(y,v)=0{,}0016$, prácticamente cero. También mejora la
equidad en calidad frente a BM y DA y mantiene un rank promedio cercano al de
WP-priority. La versión escalable de WP mejora frente a DA, pero no supera a
SED-lex.

La lectura central es que, en Bogotá, la prioridad importa más que la forma fina
del objetivo. Cuando las preferencias ya están restringidas por localidad, una
prioridad lexicográfica por vulnerabilidad con desempate geográfico redistribuye
el acceso dentro de los conjuntos de opciones locales y corrige casi por completo
el sorting ingreso--seguridad. WP sigue siendo útil como ejercicio de diseño:
muestra cómo construir prioridades estables orientadas a una métrica de equidad
específica. Pero los datos reales sugieren que el sistema vigente ya contiene una
palanca potente de corrección visual.

\paragraph{Limitaciones.}
(i)~La demanda SED no está desagregada por localidad de origen del solicitante;
por tanto, las cuotas de mercado se construyen suponiendo que la demanda de cada
colegio proviene de su localidad. (ii)~La estimación BLP opera con 19~mercados,
lo que limita las propiedades asintóticas de los errores estándar agrupados.
(iii)~El sorting residencial puede sesgar la penalización estimada por distancia:
familias de mayor ingreso pueden vivir cerca de colegios de mayor calidad y
mejores entornos urbanos. (iv)~En datos reales aplicamos WP-priority, no el LP
global de WP, por razones de escala computacional; por tanto, los resultados
reales evalúan la prioridad aprendida, no todo el politopo estable.

\paragraph{Implicaciones de política.}
El resultado de política no es que el mecanismo vigente falle, sino que cumple
una función correctiva importante. Las señales visuales influyen en la demanda,
pero SED-lex neutraliza casi por completo su transmisión a la asignación final.
Por tanto, una reforma orientada a equidad debería preservar la prioridad por
vulnerabilidad y el desempate geográfico. Cualquier intervención sobre apariencia
o infraestructura visual debería complementarse con información activa y
verificable sobre calidad académica \citep{HastingsWeinstein2008,Allende2019},
para que las preferencias familiares no dependan únicamente de señales visibles.


% ============================================================
\bibliographystyle{plainnat}

\begin{thebibliography}{99}

\bibitem[Abdulkadiro\u{g}lu y S\"onmez(2003)]{AbdulSonmez2003}
Abdulkadiro\u{g}lu, A.\ y S\"onmez, T.
\newblock School choice: A mechanism design approach.
\newblock \emph{American Economic Review}, 93(3):729--747, 2003.

\bibitem[Allende et~al.(2019)]{Allende2019}
Allende, C., Gallego, F.~A., y Neilson, C.~A.
\newblock Approximating the equilibrium effects of informed school choice.
\newblock \emph{The Review of Economic Studies}, revise and resubmit, 2019.

\bibitem[Berry(1994)]{Berry1994}
Berry, S.
\newblock Estimating discrete-choice models of product differentiation.
\newblock \emph{RAND Journal of Economics}, 25(2):242--262, 1994.

\bibitem[Conlon y Gortmaker(2020)]{ConlonGortmaker2020}
Conlon, C.\ y Gortmaker, J.
\newblock Best practices for differentiated products demand estimation with {PyBLP}.
\newblock \emph{RAND Journal of Economics}, 51(4):1108--1161, 2020.

\bibitem[DANE(2021)]{DANE2021EM}
Departamento Administrativo Nacional de Estadística~(DANE).
\newblock \emph{Encuesta Multipropósito Bogotá~2021 (EM2021)}.
\newblock DANE -- Secretaría Distrital de Planeación, 2021.

\bibitem[DANE(2024a)]{DANE2024pobreza}
Departamento Administrativo Nacional de Estadística~(DANE).
\newblock \emph{Boletín técnico: Pobreza Monetaria Colombia 2024}.
\newblock DANE, 2024.

\bibitem[DANE(2024b)]{DANE2024clases}
Departamento Administrativo Nacional de Estadística~(DANE).
\newblock \emph{Comunicado de prensa: Clases Sociales 2024}.
\newblock DANE, 2024.

\bibitem[Radford et~al.(2021)]{Radford2021}
Radford, A., Kim, J.~W., Hallacy, C., et~al.
\newblock Learning transferable visual models from natural language supervision.
\newblock En \emph{ICML~2021}, pp.~8748--8763. PMLR, 2021.

\bibitem[Dubey et~al.(2016)]{Dubey2016}
Dubey, A., Naik, N., Parikh, D., Raskar, R., y Hidalgo, C.~A.
\newblock Deep learning the city: Quantifying urban perception at a global scale.
\newblock En \emph{ECCV~2016}, pp.~196--212. Springer, 2016.

\bibitem[Gale y Shapley(1962)]{GaleShapley1962}
Gale, D.\ y Shapley, L.~S.
\newblock College admissions and the stability of marriage.
\newblock \emph{The American Mathematical Monthly}, 69(1):9--15, 1962.

\bibitem[Hastings et~al.(2009)]{Hastings2009}
Hastings, J.~S., Kane, T.~J., y Staiger, D.~O.
\newblock Heterogeneous preferences and the efficacy of public school choice
(NBER~WP~12145).
\newblock National Bureau of Economic Research, 2009.

\bibitem[Hastings y Weinstein(2008)]{HastingsWeinstein2008}
Hastings, J.~S.\ y Weinstein, J.~M.
\newblock Information, school choice, and academic achievement.
\newblock \emph{Quarterly Journal of Economics}, 123(4):1373--1414, 2008.

\bibitem[Lee y Seung(2001)]{LeeSeung2001}
Lee, D.~D.\ y Seung, H.~S.
\newblock Algorithms for non-negative matrix factorization.
\newblock En \emph{NeurIPS~2000}, pp.~556--562. MIT Press, 2001.

\bibitem[Naik et~al.(2017)]{Naik2017}
Naik, N., Raskar, R., y Hidalgo, C.~A.
\newblock Computer vision uncovers predictors of physical urban change.
\newblock \emph{PNAS}, 114(29):7571--7576, 2017.

\bibitem[Roth(1984)]{Roth1984}
Roth, A.~E.
\newblock The evolution of the labor market for medical interns and residents:
a case study in game theory.
\newblock \emph{Journal of Political Economy}, 92(6):991--1016, 1984.

\bibitem[SED(2023)]{SED2023caracterizacion}
Secretaría de Educación del Distrito~(SED).
\newblock \emph{I~Informe Caracterización Sector Educativo Bogotá D.C.\ 2023}.
\newblock Alcaldía Mayor de Bogotá, 2023.

\bibitem[SED(2025)]{Resolucion1587}
Secretaría de Educación del Distrito~(SED).
\newblock \emph{Resolución~1587 de 2025: gestión de la cobertura educativa
  para las vigencias 2025 y~2026}.
\newblock SED, 2025.

\bibitem[Narasimhan et~al.(2016)]{NarasimhanAgarwalParkes2016}
Narasimhan, H., Agarwal, S., y Parkes, D.~C.
\newblock Automated mechanism design without money via machine learning.
\newblock En \emph{IJCAI~2016}, pp.~433--439. AAAI Press, 2016.

\bibitem[Simonyan y Zisserman(2015)]{Simonyan2015}
Simonyan, K.\ y Zisserman, A.
\newblock Very deep convolutional networks for large-scale image recognition.
\newblock En \emph{ICLR~2015}, 2015.

\bibitem[Vande Vate(1989)]{VandeVate1989}
Vande Vate, J.~H.
\newblock Linear programming brings marital bliss.
\newblock \emph{Operations Research Letters}, 8(3):147--153, 1989.

\end{thebibliography}

\end{document}
